% ===================================================================
%  TABEL Nr. 11 — De stad en haar gedenktekens. — De brand.
%  Traduction 2026 d'apres texto/io, controlee sur texto/fr et
%  texto/es. Consigne : voir texto/nl/10-tabel-01.tex.
% ===================================================================

%%K t11-apar-1 apar
\VUsaut{2.0mm}
\VUcentre{13.2pt}{90}{DERDE REEKS}
\VUsaut{3.0mm}
\VUfilet{16mm}
\VUsaut{5.6mm}
\VUtitre{15.0pt}{60}{TABEL Nr. 11}
\VUsaut{1.4mm}
\VUpk{11.4pt}{40}{De stad en haar gedenktekens. --- De brand.}
\VUsaut{2.2mm}

%%K t11-tit-1 sub
\VUpk{10.2pt}{40}{De voorstad.}
\VUsaut{3.0mm}

%%K t11-c1-01-1 p
1. --- Ik woon in een grote stad die op 100 kilometer van de Oceaan ligt, en die
toch een \VUgras{haven} \textsuperscript{(1)} van de eerste rang is. De stroom
die haar begrenst heeft een breedte van 400 meter en vormt een zeer mooie rede.

%%K t11-c1-02-1 p
2. --- Het hele deel van de stad dat aan de \VUgras{kaden}
\textsuperscript{(2)} ligt ziet er geheel maritiem en industrieel uit, en vormt
een \VUgras{voorstad} \textsuperscript{(3)} die alleen door zeelieden en
arbeiders bewoond wordt. Dezen werken in de \VUgras{fabrieken}
\textsuperscript{(4)} en in de \VUgras{gasfabriek} \textsuperscript{(5)} waarvan
de hoge \VUgras{schoorsteen} \textsuperscript{(6)} en de \VUgras{gashouder} van
zeer ver te zien zijn.

%%K t11-c1-03-1 p
3. --- In de middeleeuwen was de stad versterkt, en men vindt nog enkele resten
van haar wallen, in het bijzonder de \VUgras{poort} \textsuperscript{(8)} van
het oude \VUgras{versterkte kasteel} \textsuperscript{(7)} geflankeerd door zijn
twee \VUgras{torens} \textsuperscript{(9)} die nu als \VUgras{gevangenis}
dienen.

%%K t11-c1-04-1 p
4. --- In die hele \VUgras{wijk}, die er zeer oud uitziet, is maar één gebouw
betrekkelijk modern : dat is het \VUgras{douanekantoor} \textsuperscript{(10)}.
Het ligt tussen de kade en het \VUgras{straatje} \textsuperscript{(11)} dat naar
het oude \VUgras{gemeentehuis} \textsuperscript{(12)} leidt. Men herkent dat
laatste gedenkteken gemakkelijk aan de reusachtige \VUgras{klokkentoren}
\textsuperscript{(13)} die de oude stad beheerst.

%%K t11-c1-05-1 p
5. --- Aan de andere kant van de straat \VUgras{staat} een gebouw dat in
dezelfde stijl als het douanekantoor gebouwd is : dat is de \VUgras{Beurs}
\textsuperscript{(14)}. Een van haar gevels kijkt uit op een klein plein waar het
\VUgras{prefectuurgebouw} \textsuperscript{(15)} staat. Onze stad is namelijk,
zonder een hoofdstad te zijn, toch een belangrijke departementshoofdstad.

%%K t11-tit-2 sub
\VUsaut{5.0mm}
\VUpk{10.2pt}{40}{Het hoogste deel van de stad.}
\VUsaut{3.4mm}

%%K t11-c1-06-1 p
6. --- Het garnizoen, tamelijk talrijk, is ondergebracht in ruime zeer gezonde
\VUgras{kazernes} \textsuperscript{(16)}; zij strekken zich uit tot aan het hek
van het \VUgras{gemeentepark} \textsuperscript{(17)}. De \VUgras{boulevard}
\textsuperscript{(19)} die naar dat park leidt loopt achter de
\VUgras{spaarbank} \textsuperscript{(18)} langs en eindigt op het plein van de
\VUgras{kathedraal} \textsuperscript{(20)}.

%%K t11-c1-07-1 p
7. --- Al die gedenktekens strekken zich als een amfitheater uit op een heuvel
waar ook ons ruime \VUgras{lyceum} \textsuperscript{(21)} staat.

%%K t11-c1-07-2 p
Als men langs een wat hellende weg afdaalt, die op de Grote straat uitkomt, komt
men bij het \VUgras{Gerechtshof} \textsuperscript{(22)} waarvoor een aangename
\VUgras{openbare tuin} \textsuperscript{(26)} zich uitstrekt. Dat
\VUgras{plantsoen} is niet erg groot, maar het vormt midden in de stad een oase
van groen en koelte. Daarom wordt het veel bezocht door de stedelingen die graag
onder het tentdak van het \VUgras{café} \textsuperscript{(25)} komen zitten, om
te luisteren naar de militaire muzikanten die op donderdag en zondag in de
\VUgras{kiosk} \textsuperscript{(27)} spelen die tussen de grasperken en de
struikgroepen staat.

%%K t11-c1-08-1 p
8. --- Op een van die grasperken staat een bronzen \VUgras{standbeeld}
\textsuperscript{(28)}, een gedenkteken opgericht ter nagedachtenis van een van
onze stadgenoten, die grote diensten aan zijn geboortestad bewees.

%%K t11-tit-3 sub
\VUsaut{4.0mm}
\VUpk{10.2pt}{40}{De vervoermiddelen.}
\VUsaut{3.4mm}

%%K t11-c1-09-1 p
9. --- Tot nu toe wordt onze stad nog niet met elektrisch licht verlicht, zij
heeft alleen \VUgras{gaslantaarns} \textsuperscript{(29)}. Maar de
\VUgras{plaatselijke pers} voert campagne opdat dat verlichtingsstelsel verbeterd
zou worden, \VUgras{zoals} men al \VUgras{het tractiestelsel van de} trams
volmaakt heeft. \VUgras{De} oude rijtuigen, \VUgras{getrokken} door
\VUgras{paarden}, werden praktisch vervangen door \VUgras{elektrische trams}
\textsuperscript{(31)} die de reizigers snel van het ene uiteinde van de
\VUgras{stad naar het andere} brengen, tegen een zeer \VUgras{lage} prijs.

%%K t11-c1-10-1 p
10. --- Bij het Gerechtshof staat de \VUgras{Bank} \textsuperscript{(32)} waar
men de handelswaarden disconteert. De \VUgras{bankbedienden}
\textsuperscript{(33)} gaan de schulden aan huis innen.

%%K t11-tit-4 sub
\VUsaut{4.0mm}
\VUpk{10.2pt}{40}{Kunst en onderwijs.}
\VUsaut{3.4mm}

%%K t11-c1-11-1 p
11. --- Het mooie gebouw dat daarnaast staat, is het \VUgras{museum}. Het bevat
tegelijk de \VUgras{bibliotheek} \textsuperscript{(34)} van de stad, de mooie
\VUgras{natuurwetenschappelijke verzamelingen} \textsuperscript{(35)}
(Natuurhistorisch museum) en een sierlijk \VUgras{schilderijenmuseum}
\textsuperscript{(36)} dat een prachtige blijvende \VUgras{tentoonstelling}
vormt.

%%K t11-c1-11-2 p
Het gaat tweemaal per week voor het publiek open, maar de vreemdelingen mogen
het elke dag bezoeken tegen een fooi aan de \VUgras{bewaker}
\textsuperscript{(37)}. De \VUgras{schilders} \textsuperscript{(38)} komen daar
werken. Men ziet hen \VUgras{voor} hun ezel, met een \VUgras{palet} in de hand,
de beroemde schilderijen kopiëren.

%%K t11-c1-12-1 p
12. --- Op de hoogte, achter het museum, beginnen wij de \VUgras{koepel}
\textsuperscript{(40)} van het \VUgras{ziekenhuis} \textsuperscript{(39)} te
zien en de gebouwen van de \VUgras{kweekschool} \textsuperscript{(41)} waarin de
onderwijzers van onze gemeentescholen opgeleid werden. Op het Theaterplein staat
een \VUgras{postkantoor} \textsuperscript{(43)} dat in een \VUgras{particulier
huis} \textsuperscript{(42)} gevestigd is.

%%K t11-tit-5 sub
\VUsaut{5.0mm}
\VUpk{11.4pt}{40}{De brand.}
\VUsaut{3.0mm}

%%K t11-tit-6 sub
\VUpk{10.2pt}{40}{Brandkreet.}
\VUsaut{3.4mm}

%%K t11-c2-01-1 p
1. --- Een verschrikkelijk gerucht verspreidt zich door de stad : er is zojuist
brand in het theater uitgebroken! Op het ogenblik dat ik op het \VUgras{plein}
\textsuperscript{(96)} kom, staat de menigte al op het \VUgras{trottoir}
\textsuperscript{(46)} en op de \VUgras{rijweg} \textsuperscript{(47)}
verzameld. De \VUgras{postbeambten} \textsuperscript{(44)} en de
\VUgras{brievenbestellers} \textsuperscript{(45)} komen het postkantoor uit. Een
\VUgras{trambestuurder} \textsuperscript{(48)} moest zijn (tram)wagen laten
stilhouden om een gemeentelijke \VUgras{ziekenwagen} \textsuperscript{(50)} te
laten passeren, die met een \VUgras{chirurg} \textsuperscript{(52)} en een
\VUgras{verpleger} \textsuperscript{(51)} komt om een \VUgras{gewonde}
\textsuperscript{(53)} te verzorgen die op een \VUgras{draagbaar}
\textsuperscript{(54)} bij de \VUgras{krantenkiosk} \textsuperscript{(55)}
gebracht is.

%%K t11-c2-02-1 p
2. --- Een compagnie infanterie, die van de oefeningen terugkeerde, hield stil om
de ordedienst te verzekeren samen met de \VUgras{bereden gemeentepolitie}
\textsuperscript{(61)}. \VUgras{De ruiters}, wier paarden steigeren, doen, beter dan de
\VUgras{soldaten} \textsuperscript{(57)} of de \VUgras{gendarmes}
\textsuperscript{(63)}, de \VUgras{nieuwsgierigen} \textsuperscript{(56)}
achteruitwijken. Overigens hebben de gendarmes het zeer druk. Een van hen
wijst de \VUgras{politieagenten} \textsuperscript{(64)} waarheen men een andere
\VUgras{gewonde} \textsuperscript{(65)} moet vervoeren, een moedige brandweerman
die van een ladder gevallen is.

%%K t11-c2-03-1 p
3. --- \VUgras{De officier} \textsuperscript{(66)} geeft nauwkeurig de plaats
aan waar men de aankomende \VUgras{stoomspuit} \textsuperscript{(67)} moet
plaatsen. In afwachting van die krachtige machine heeft hij een
\VUgras{handspuit} \textsuperscript{(68)} in stelling gebracht waarvan de lange
\VUgras{slang} \textsuperscript{(69)} het water als een stortvloed op het vuur
werpt. Zes brandweerlieden bedienen haar, terwijl hun kameraad, die de
\VUgras{straalpijp} \textsuperscript{(70)} vasthoudt, de \VUgras{straal}
\textsuperscript{(71)} naar het middelpunt van de brand richt.

%%K t11-c2-04-1 p
4. --- Aan de andere kant van het theater heeft men de grote \VUgras{ladder}
\textsuperscript{(72)} geplaatst. Zij stelt de brandweerlieden in staat de
vlammen te naderen die tussen wervelingen van zwarte \VUgras{rook}
\textsuperscript{(75)} opspringen.

%%K t11-tit-7 sub
\VUsaut{5.0mm}
\VUpk{10.2pt}{40}{Dappere daden.}
\VUsaut{3.4mm}

%%K t11-c2-05-1 p
5. --- \VUgras{De brand} \textsuperscript{(74)} is ontstaan in de foyer van het
\VUgras{theater} \textsuperscript{(73)}, terwijl men \VUgras{de opera}
repeteerde waarvan de eerste \VUgras{voorstelling} morgen zou hebben
plaatsgevonden. Gelukkig waren er maar weinig \VUgras{toeschouwers}
\textsuperscript{(76)} in de zaal. De ongelukkigen zijn op het \VUgras{balkon}
\textsuperscript{(77)} gevlucht waar moedige \VUgras{brandweerlieden}
\textsuperscript{(86)} hun met ladders en touwen te hulp komen.

%%K t11-c2-06-1 p
6. --- Onder het \VUgras{peristilium} \textsuperscript{(78)}, waar het
\VUgras{aanplakbiljet} \textsuperscript{(79)} over de voorstelling inlicht, ziet
men de \VUgras{muzikanten} \textsuperscript{(81)} van het \VUgras{orkest}
\textsuperscript{(80)} vluchten, hun instrumenten meedragend : \VUgras{viool}
\textsuperscript{(81)}, \VUgras{fluit} \textsuperscript{(82)},
\VUgras{violoncel} \textsuperscript{(83)}, \VUgras{trombone}
\textsuperscript{(84)}, \VUgras{trom} \textsuperscript{(85)}, enz.. Men tracht
een deel van het \VUgras{meubilair} \textsuperscript{(87)} te redden.

%%K t11-c2-07-1 p
7. --- \VUgras{De acteurs} \textsuperscript{(89)} en \VUgras{de actrices}
\textsuperscript{(88)}, \VUgras{zangers} en \VUgras{zangeressen} moesten in hun
\VUgras{toneelkostuum} \textsuperscript{(90)} vluchten. De tenor legt aan een
\VUgras{journalist} \textsuperscript{(92)} uit hoe het gebeurd is. \VUgras{De
verslaggever} kwam dadelijk vanaf het eerste ogenblik toegelopen met de
overheidspersonen : de \VUgras{prefect} \textsuperscript{(91)}, de
\VUgras{burgemeester} \textsuperscript{(93)}, de \VUgras{politiecommissaris}
\textsuperscript{(94)} en de \VUgras{rechterlijke ambtenaar}
\textsuperscript{(95)}, die een onderzoek zullen instellen om over de
verantwoordelijkheden te oordelen.
\VUsaut{6.0mm}
\VUfilet{22mm}
